
\[
  0\longrightarrow \Projz(\underline{L}_{r+1})\longrightarrow\ldots\longrightarrow
  \Projz(\underline{L}_0)\longrightarrow\Projz(\underline{M})\longrightarrow 0,
\]
et l'égalité suivante dans \(K^{\bullet}(X)\) :
\[
  \cl(\underline{F})=\sum_i(-1)^i\cl\bigl(\Projz(\underline{L}_i)\bigr).
\]
En utilisant la filtration sur les \(\underline{L}_i\) définie en 1.7, on trouve :
\[
  \cl(\Projz(\underline{L}_i))=\sum_j \cl\bigl(\Projz(\underline{N}_{ij}\otimes_{\underline{O}_S}\underline{B}(j))\bigr)
\]
Or \(\Projz(\underline{N}_{ij}\otimes_{\underline{O}_S}\underline{B}(j))=f^*(\underline{N}_{ij})(j)\)~; on trouve donc :
\[
\begin{aligned}
  \cl(\Projz(\underline{L}_i))
  &=\sum_j \cl\bigl(f^*(\underline{N}_{ij})\bigr).\xi^j \\
  &=\sum_j f^*\bigl(\cl(\underline{N}_{ij})\bigr).\xi^j.
\end{aligned}
\]
On a donc finalement exprimé \(\cl(\underline{F})\) comme combinaison linéaire à coefficients dans \(K^{\bullet}(S)\) des puissances de \(\xi\).

1.9. Pour voir que \(1,\xi,\ldots,\xi^{r+1}\) sont linéairement dépendants sur \(K^{\bullet}(S)\) on considère l'homomorphisme canonique surjectif : \(f^*(\underline{E})\to\underline{O}_X(1)\), dont le noyau \(\underline{F}\) est localement libre de rang \(r\). On peut donc écrire :
\[
  \lambda^{r+1}\bigl(f^*(E)-\xi\bigr)=0,
  \qquad \text{avec } E=\cl\underline{E},
\]
les opérations \(\lambda^i\) sur \(K^{\bullet}(X)\) étant définies dans V 2.2. On obtient donc :
\[
  \sum_{i=0}^{r+1}\lambda^i(f^*(E))\lambda^{r+1-i}(-\xi)=0.
\]
Comme \(\underline{O}_X(1)\) est un faisceau inversible, \(\lambda_t(\xi)=1+\xi t\), et par suite, \(\lambda_t(-\xi)=1/(1+\xi t)\), d'où finalement la relation :
\begin{equation}\tag{1.9.1}
  \sum_{i=0}^{r+1}(-1)^{r+1-i}\xi^{r+1-i}\lambda^i(E)=0.
\end{equation}

Montrons enfin que \(1,\xi,\ldots,\xi^r\) sont indépendants sur \(K^{\bullet}(S)\)\SGAAuthorFootnote{*}{Cette démonstration est due à J.P. Jouanolou.}.


